% GENERATED FILE. Edit canonical lessons, then run npm run generate:books.
% canonical-source-hash: fnv1a64:d85c5c2e01671146
% canonical-lessons: LA-C54-edo, LA-C54-bibo, LA-C54-oculus, LA-C54-auris, LA-C54-pes

\chapter{Edo, Bibo, Oculus, Auris, Pes}
\label{ch:la-corpus}
\hlchaptermodality{54}

\begin{chapteropening}
\textbf{By the end of this chapter:} \emph{I can say that I eat and drink, and name the eye, the ear and the foot beside the head and hand I already had.}

The fifth word of the group, taken apart against the four before it.
\end{chapteropening}

\section[edō, esse]{edō — and the infinitive that collides with to be}

\label{lesson:LA-C54-edo}



\begin{quote}
Six recalls before the new word, at growing distances. They are not decoration: each one is the retrieval window this book measures, and each names a word far enough back that saying it is work.

\begin{itemize}
\raggedright
  \item \emph{Recall:} say \emph{liber} — \textbf{R1}, one lesson back, before it decays
  \item \emph{Recall:} say \emph{iterum} — \textbf{R2}, the first real retrieval, five to fifteen lessons back
  \item \emph{Recall:} say \emph{hic} — \textbf{R3}, durable, twenty to sixty lessons back
  \item \emph{Recall:} say \emph{currō} — \textbf{R3}, durable again, a second word from the same distance
  \item \emph{Recall:} say \emph{bene} — \textbf{R4}, recognition at distance, eighty lessons back or more
  \item \emph{Recall:} say \emph{diēs} — \textbf{R4}, recognition at distance, a second word from that far back
\end{itemize}
\end{quote}

\begin{cousinweb}
\textbf{edō} — \textquotedblleft{}\textbf{I eat}.\textquotedblright{}

Fifty-three chapters gave you bread, salt, honey, milk, eggs, fish, water and wine, and no verb to do anything with them. Here it is.

The root is Proto-Indo-European \emph{\textbf{*h\textsubscript{1}ed-}}, and English \textbf{eat} is an \textbf{inherited cousin}, not a loan — the same relationship \emph{piscis} and \emph{fish} have, by the same Grimm's Law. So are \textbf{etch} and, at a distance, \textbf{fret} (\emph{fra-etan}, \textquotedblleft{}to eat away\textquotedblright{}).

And then a trap worth meeting once. The infinitive of \emph{edō} is \textbf{ēsse}, and the infinitive of \emph{sum}, \textquotedblleft{}to be,\textquotedblright{} is \textbf{esse}. Four identical letters again. Only the long \emph{ē} separates eating from being, and only in speech; a Roman reader took it from the sentence.

\textbf{Edible} and \textbf{comestible} are loans from this verb.
\end{cousinweb}

\subsection*{Your turn}
Say these aloud:

\begin{itemize}
\raggedright
  \item \emph{edō, esse}
  \item every word this chapter has given you so far, in order
\end{itemize}

\begin{tcolorbox}[breakable,colback=teal!4,colframe=teal!35!black,title={Before you move on}]
Is English \emph{eat} a cousin or a loan? (\textbf{A cousin}.) What two verbs share the spelling \emph{esse}? (\emph{\textbf{Edō}} and \emph{\textbf{sum}} — separated by vowel length.)
\end{tcolorbox}
\section[bibō, bibere]{bibō — the verb inside beverage, bib and imbibe}

\label{lesson:LA-C54-bibo}



\begin{quote}
Six recalls before the new word, at growing distances. They are not decoration: each one is the retrieval window this book measures, and each names a word far enough back that saying it is work.

\begin{itemize}
\raggedright
  \item \emph{Recall:} say \emph{edō} — \textbf{R1}, one lesson back, before it decays
  \item \emph{Recall:} say \emph{magister} — \textbf{R2}, the first real retrieval, five to fifteen lessons back
  \item \emph{Recall:} say \emph{iste} — \textbf{R3}, durable, twenty to sixty lessons back
  \item \emph{Recall:} say \emph{aperiō} — \textbf{R3}, durable again, a second word from the same distance
  \item \emph{Recall:} say \emph{diēs} — \textbf{R4}, recognition at distance, eighty lessons back or more
  \item \emph{Recall:} say \emph{nox} — \textbf{R4}, recognition at distance, a second word from that far back
\end{itemize}
\end{quote}

\begin{cousinweb}
\textbf{bibō} — \textquotedblleft{}\textbf{I drink}.\textquotedblright{}

The other half of the pair, and the one that finally does something with the \emph{aqua} and \emph{vīnum} of the bread-and-wine lesson.

\begin{quote}
\textbf{Aquam bibō.} — I drink water. \textbf{Aquam aut vīnum bibō?} — Do I drink water or wine? (with the \emph{aut} of the denying chapter.)
\end{quote}

The form is odd and the oddity is old: \emph{bibō} is a \textbf{reduplicated} verb, the root syllable said twice, from Proto-Indo-European \emph{\textbf{*peh\textsubscript{3}-}}. Latin regularly turns that initial \emph{p} into \emph{b} here, which is why \emph{bibō} looks nothing like Greek \emph{pīnō} although they are the same word.

English \textbf{beverage} is Old French \emph{bevrage}, from \emph{bibere}; \textbf{imbibe} and \textbf{bibulous} are direct loans; and a baby's \textbf{bib} is probably from the same verb, the thing you drink over — probably, not certainly.
\end{cousinweb}

\subsection*{Your turn}
Say these aloud:

\begin{itemize}
\raggedright
  \item \emph{bibō, bibere}
  \item every word this chapter has given you so far, in order
\end{itemize}

\begin{tcolorbox}[breakable,colback=teal!4,colframe=teal!35!black,title={Before you move on}]
Say \textquotedblleft{}I drink water.\textquotedblright{} (\emph{\textbf{Aquam bibō}}.) Is \emph{bib} certainly from \emph{bibō}? (\textbf{Probably, not certainly}.)
\end{tcolorbox}
\section[oculus, oculī]{oculus — and the eye already hiding inside a word you know}

\label{lesson:LA-C54-oculus}



\begin{quote}
Six recalls before the new word, at growing distances. They are not decoration: each one is the retrieval window this book measures, and each names a word far enough back that saying it is work.

\begin{itemize}
\raggedright
  \item \emph{Recall:} say \emph{bibō} — \textbf{R1}, one lesson back, before it decays
  \item \emph{Recall:} say \emph{discipulus} — \textbf{R2}, the first real retrieval, five to fifteen lessons back
  \item \emph{Recall:} say \emph{ille} — \textbf{R3}, durable, twenty to sixty lessons back
  \item \emph{Recall:} say \emph{claudō} — \textbf{R3}, durable again, a second word from the same distance
  \item \emph{Recall:} say \emph{nox} — \textbf{R4}, recognition at distance, eighty lessons back or more
  \item \emph{Recall:} say \emph{salvē / salvēte} — \textbf{R4}, recognition at distance, a second word from that far back
\end{itemize}
\end{quote}

\begin{cousinweb}
\textbf{oculus} — \textquotedblleft{}\textbf{an eye}.\textquotedblright{}

The third body word this book has, after \emph{caput} and \emph{manus}.

English \textbf{eye} is an \textbf{inherited cousin} — Proto-Indo-European \emph{\textbf{*h\textsubscript{3}ek\textsuperscript{w}-}} gave Latin \emph{oculus} and Germanic \emph{augon}, which is \emph{eye}. And the loans are everywhere: \textbf{ocular}, \textbf{binoculars}, \textbf{monocle}, \textbf{oculist}.

The best of them is \textbf{inoculate}. Latin \emph{inoculāre} was a \textbf{gardening} word: to graft a bud — an \emph{oculus}, an \textquotedblleft{}eye\textquotedblright{} — into a plant. Eighteenth-century physicians borrowed the gardener's word for grafting a little disease into a person, and that is why an injection is named after an eye.

Italian \emph{occhio} and Spanish \emph{ojo} both descend from it, by the same \emph{-cl-} to \emph{-cchi-} change Italian makes on \emph{occhio}.
\end{cousinweb}

\subsection*{Your turn}
Say these aloud:

\begin{itemize}
\raggedright
  \item \emph{oculus, oculī}
  \item every word this chapter has given you so far, in order
\end{itemize}

\begin{tcolorbox}[breakable,colback=teal!4,colframe=teal!35!black,title={Before you move on}]
Is English \emph{eye} a cousin or a loan? (\textbf{A cousin}.) Why is \emph{inoculate} named after an eye? (\textbf{It was a gardening word} — grafting a bud.)
\end{tcolorbox}
\section[auris, auris]{auris — and the reason a Roman listened by obeying}

\label{lesson:LA-C54-auris}



\begin{quote}
Six recalls before the new word, at growing distances. They are not decoration: each one is the retrieval window this book measures, and each names a word far enough back that saying it is work.

\begin{itemize}
\raggedright
  \item \emph{Recall:} say \emph{oculus} — \textbf{R1}, one lesson back, before it decays
  \item \emph{Recall:} say \emph{schola} — \textbf{R2}, the first real retrieval, five to fifteen lessons back
  \item \emph{Recall:} say \emph{is} — \textbf{R3}, durable, twenty to sixty lessons back
  \item \emph{Recall:} say \emph{ferō} — \textbf{R3}, durable again, a second word from the same distance
  \item \emph{Recall:} say \emph{salvē / salvēte} — \textbf{R4}, recognition at distance, eighty lessons back or more
  \item \emph{Recall:} say \emph{avē / avēte} — \textbf{R4}, recognition at distance, a second word from that far back
\end{itemize}
\end{quote}

\begin{cousinweb}
\textbf{auris} — \textquotedblleft{}\textbf{an ear}.\textquotedblright{}

English \textbf{ear} is an \textbf{inherited cousin} again — the same Proto-Indo-European \emph{\textbf{*h\textsubscript{2}ews-}} on both sides.

And it connects to a verb you already have. The \emph{audiō} lesson gave you \textquotedblleft{}I hear,\textquotedblright{} and noted that \textbf{obey} is hidden inside it. Now you can see why: \emph{audiō} is itself built on this root — \emph{aus-diō}, roughly \textquotedblleft{}I put the ear to\textquotedblright{} — so hearing, in Latin, is an action performed with the named organ. \emph{Ob-audīre}, \textquotedblleft{}to listen towards,\textquotedblright{} wore down through French into \textbf{obey}.

\textbf{Auricle}, \textbf{aural} and \textbf{auscultate} are loans. Watch the near-collision: \textbf{aural} (of the ear) and \textbf{oral} (of the mouth) are different words from different roots, and English pronounces them identically — which is why careful writers avoid both.
\end{cousinweb}

\subsection*{Your turn}
Say these aloud:

\begin{itemize}
\raggedright
  \item \emph{auris, auris}
  \item every word this chapter has given you so far, in order
\end{itemize}

\begin{tcolorbox}[breakable,colback=teal!4,colframe=teal!35!black,title={Before you move on}]
Which earlier verb is built on this root? (\emph{\textbf{Audiō}}.) What is the difference between \emph{aural} and \emph{oral}? (\textbf{Ear against mouth} — different roots, same sound.)
\end{tcolorbox}
\section[pēs, pedis]{pēs — the clearest Grimm's Law pair in the book, and the last word of the tranche}

\label{lesson:LA-C54-pes}



\begin{quote}
Six recalls before the new word, at growing distances. They are not decoration: each one is the retrieval window this book measures, and each names a word far enough back that saying it is work.

\begin{itemize}
\raggedright
  \item \emph{Recall:} say \emph{auris} — \textbf{R1}, one lesson back, before it decays
  \item \emph{Recall:} say \emph{liber} — \textbf{R2}, the first real retrieval, five to fifteen lessons back
  \item \emph{Recall:} say \emph{ibi} — \textbf{R3}, durable, twenty to sixty lessons back
  \item \emph{Recall:} say \emph{accipiō} — \textbf{R3}, durable again, a second word from the same distance
  \item \emph{Recall:} say \emph{videō} — \textbf{R4}, recognition at distance, eighty lessons back or more
  \item \emph{Recall:} say \emph{salvē / salvēte} — \textbf{R4}, recognition at distance, a second word from that far back
\end{itemize}
\end{quote}

\begin{cousinweb}
\textbf{pēs, pedis} — \textquotedblleft{}\textbf{a foot}.\textquotedblright{}

The \emph{piscis} lesson taught Grimm's Law on \emph{piscis} and \textbf{fish}, and promised the other pairs. This is the cleanest of them: Latin \emph{p} against English \emph{f}, \emph{pēs} against \textbf{foot}, on the same Proto-Indo-European \emph{\textbf{*ped-}}.

And the genitive matters again, exactly as it did for \emph{mīles}. The nominative is \emph{pēs}; the stem is \emph{ped-}; and every English loan is built on the stem: \textbf{pedal}, \textbf{pedestrian}, \textbf{impede} (\textquotedblleft{}to shackle the feet\textquotedblright{}), \textbf{expedite} (\textquotedblleft{}to free them\textquotedblright{}), \textbf{biped}, \textbf{pedicure}.

\textbf{Pedigree} is the best of them: French \emph{pied de grue}, \textquotedblleft{}crane's foot,\textquotedblright{} because the three-line mark that shows descent on a genealogical chart looks like a bird's footprint.

That is thirty-five words in seven chapters. Say the five of this chapter in order, then say the head and the hand you have had from the start, and you have named a body.
\end{cousinweb}

\subsection*{Your turn}
Say these aloud:

\begin{itemize}
\raggedright
  \item \emph{pēs, pedis}
  \item every word this chapter has given you so far, in order
\end{itemize}

\begin{tcolorbox}[breakable,colback=teal!4,colframe=teal!35!black,title={Before you move on}]
What sound does Latin \emph{p} correspond to in inherited English? (\textbf{F} — Grimm's Law.) Which form do the English loans use? (\textbf{The genitive stem} \emph{ped-}.) And why is a \emph{pedigree} a foot? (\textbf{The chart mark looks like a crane's footprint}.)
\end{tcolorbox}
